% ============================================================
%  GUÍA DE REPASO — TEMA 1 (PRIMER LAPSO)
%  Números Racionales (Q): Operaciones, Propiedades,
%  Operaciones Combinadas, Potenciación y Ecuaciones
%  Curso: 2do Año — CELAM
%  Autor: Ing. Gabriel Astudillo
%  Basada en las clases del 28/09/26 (memoria_clases.md)
% ============================================================

\documentclass[12pt, letterpaper]{article}

% ── Paquetes ─────────────────────────────────────────────────

\usepackage{mathwizards-guia}
\mwguia{Guía de Repaso — Tema 1: Números Racionales ($\mathbb{Q}$)}{Ing. Gabriel Astudillo $\cdot$ 2do Año}

% ── Comandos propios de esta guía ────────────────────────────
\newcommand{\Q}{\mathbb{Q}}
\newcommand{\Z}{\mathbb{Z}}
\newcommand{\N}{\mathbb{N}}

\begin{document}
% ============================================================

% ── Portada / Título ─────────────────────────────────────────
\mwportada{Guía de Repaso — Tema 1}{Números Racionales ($\Q$)}{Operaciones, Propiedades, Potenciación y Ecuaciones}

\vspace{4pt}
\begin{cuadroinfo}
  \small
  \textbf{Presentación:} ¡Bienvenido a tu guía de repaso! Aquí reunimos \emph{todo} lo visto en el \textbf{Tema 1 del Primer Lapso}: el conjunto de los números racionales ($\Q$), las fracciones (propias, impropias y mixtas), las operaciones básicas, las propiedades de la suma y la multiplicación, las operaciones combinadas, la potenciación en $\Q$ y las ecuaciones lineales. Cada sección empieza con la \textbf{teoría} resumida, sigue con \textbf{ejemplos resueltos} paso a paso y termina con \textbf{ejercicios progresivos} de lo más sencillo a los verdaderos retos. Al final está la \textbf{clave de respuestas}.
  \\[4pt]
  \textbf{Nivel de Dificultad:}\quad
  \facil{} Básico / Directo \quad
  \medio{} Intermedio \quad
  \dificil{} Desafiante \quad
  \muyDificil{} Reto
\end{cuadroinfo}

\vspace{8pt}

% ============================================================
\section{El Conjunto de los Números Racionales ($\Q$)}
% ============================================================

\subsection{\textquestiondown Qué es un número racional?}

Los números que se pueden escribir como el \emph{cociente de dos enteros} forman el conjunto de los \textbf{números racionales}.

\begin{cuadroteoria}
  \textbf{Definición:}
  \[
    \Q = \left\{\ \frac{a}{b} \ \middle|\ a, b \in \Z,\ b \neq 0 \right\}
  \]
  donde:
  \begin{itemize}[leftmargin=*]
    \item $a$ es el \textbf{numerador}: cuántas partes se toman.
    \item $b$ es el \textbf{denominador}: en cuántas partes iguales se divide el todo ($b \neq 0$).
  \end{itemize}
  Todo número entero es racional, porque $5 = \dfrac{5}{1}$. Así:
  \[
    \N \subset \Z \subset \Q
  \]
\end{cuadroteoria}

\subsection{Tipos de fracciones y sus conversiones}

\begin{cuadroteoria}
  \begin{itemize}[leftmargin=*]
    \item \textbf{Fracción propia:} el numerador es \emph{menor} que el denominador ($\frac{3}{5}$). Representa menos de un entero.
    \item \textbf{Fracción impropia:} el numerador es \emph{mayor o igual} que el denominador ($\frac{7}{4}$). Representa más de un entero.
    \item \textbf{Número mixto:} combina un entero con una fracción propia ($2\frac{1}{3}$).
    \item \textbf{Fracciones equivalentes:} representan la misma cantidad ($\frac{1}{2} = \frac{2}{4} = \frac{3}{6}$).
  \end{itemize}
\end{cuadroteoria}

\begin{cuadropasos}
  \textbf{Convertir mixta a impropia:} $a\dfrac{b}{c} = \dfrac{a \cdot c + b}{c}$ \quad (entero por denominador, más numerador).

  \textbf{Convertir impropia a mixta:} divide el numerador entre el denominador; el \textbf{cociente} es el entero, el \textbf{resto} es el nuevo numerador y el divisor se conserva.

  \textbf{Simplificar (fracción irreductible):} divide numerador y denominador por un divisor común hasta que no queden divisores comunes distintos de $1$.
\end{cuadropasos}

\noindent\textbf{Ejemplo resuelto 1:} Clasifica $\dfrac{2}{7}$, $\dfrac{9}{4}$, $\dfrac{5}{8}$, $\dfrac{11}{3}$ y escribe las impropias como mixtas.

Las propias son $\frac{2}{7}$ y $\frac{5}{8}$ ($2<7$ y $5<8$). Para las impropias: $9 \div 4 = 2$ con resto $1$, así $\dfrac{9}{4} = 2\dfrac{1}{4}$; \quad $11 \div 3 = 3$ con resto $2$, así $\dfrac{11}{3} = 3\dfrac{2}{3}$.

\noindent\textbf{Ejemplo resuelto 2:} Escribe $3\dfrac{2}{5}$ como fracción impropia.

$$3\frac{2}{5} = \frac{3 \cdot 5 + 2}{5} = \frac{17}{5}$$

\subsection{Ejercicios: Fracciones y Conversiones}

\begin{enumerate}[leftmargin=*]
  \item \facil{} Clasifica en propia o impropia: $a)\; \dfrac{5}{9} \qquad b)\; \dfrac{7}{3} \qquad c)\; \dfrac{11}{11} \qquad d)\; \dfrac{3}{10}$

  \item \facil{} Escribe como fracción impropia: $a)\; 2\dfrac{3}{4} \qquad b)\; 5\dfrac{1}{6} \qquad c)\; 8\dfrac{2}{3}$

  \item \facil{} Escribe como número mixto: $a)\; \dfrac{17}{5} \qquad b)\; \dfrac{22}{7} \qquad c)\; \dfrac{31}{4}$

  \item \facil{} Simplifica hasta la fracción irreductible: $a)\; \dfrac{8}{12} \qquad b)\; \dfrac{24}{36} \qquad c)\; \dfrac{50}{75}$

  \item \facil{} Completa para obtener fracciones equivalentes: $a)\; \dfrac{2}{5} = \dfrac{?}{15} \qquad b)\; \dfrac{7}{8} = \dfrac{?}{32} \qquad c)\; \dfrac{5}{6} = \dfrac{20}{?}$

  \item \medio{} Ordena de menor a mayor: $\dfrac{3}{4},\ \dfrac{5}{6},\ \dfrac{7}{12}$. (Pista: lleva todo a denominador común $12$.)

  \item \medio{} Una pizza se divide en $8$ porciones iguales y me como $3$. \textquestiondown Qué fracción me queda? \textquestiondown Es propia o impropia?

  \item \dificil{} Un atleta corre $\dfrac{3}{4}$ de kilómetro cada día. \textquestiondown Qué distancia recorre en $2$ días? Expresa el resultado como mixta y como impropia.

  \item \dificil{} Escribe una fracción impropia equivalente a $5\dfrac{3}{8}$ y luego regrésala a número mixto para comprobar.
\end{enumerate}

\newpage

% ============================================================
\section{Operaciones Básicas con Fracciones}
% ============================================================

\subsection{Suma y resta de fracciones}

\begin{cuadroteoria}
  \textbf{Con igual denominador:} se suman (o restan) los numeradores y se conserva el denominador.
  \[
    \frac{a}{c} + \frac{b}{c} = \frac{a + b}{c}, \qquad \frac{a}{c} - \frac{b}{c} = \frac{a - b}{c}
  \]

  \textbf{Producto cruzado:} para $\dfrac{a}{b} \pm \dfrac{c}{d}$:
  \[
    \frac{a}{b} \pm \frac{c}{d} = \frac{a \cdot d \pm c \cdot b}{b \cdot d}
  \]

  \textbf{Método del mcm:} se calcula el mínimo común múltiplo de los denominadores, se divide el mcm entre cada denominador y se multiplica por su numerador.
\end{cuadroteoria}

\begin{cuadroadvertencia}
  \textbf{Error clásico:} \textquestiondown sumar numeradores \emph{y} denominadores? \textquestiondown No!
  \[
    \frac{1}{2} + \frac{1}{3} \neq \frac{2}{5} \qquad \text{el resultado correcto es } \frac{5}{6}
  \]
  Recuerda además \textbf{simplificar} la respuesta final.
\end{cuadroadvertencia}

\noindent\textbf{Ejemplo resuelto 1:} $\dfrac{2}{7} + \dfrac{3}{7} = \dfrac{5}{7}$ \qquad y \qquad $\dfrac{5}{9} - \dfrac{2}{9} = \dfrac{3}{9} = \dfrac{1}{3}$.

\noindent\textbf{Ejemplo resuelto 2 (producto cruzado):} $\dfrac{2}{3} + \dfrac{1}{4} = \dfrac{2 \cdot 4 + 1 \cdot 3}{3 \cdot 4} = \dfrac{8 + 3}{12} = \dfrac{11}{12}$.

\noindent\textbf{Ejemplo resuelto 3 (mcm):} $\dfrac{5}{6} - \dfrac{3}{4}$. Como mcm$(6,4) = 12$: $\dfrac{10}{12} - \dfrac{9}{12} = \dfrac{1}{12}$.

\subsection{Ejercicios: Suma y Resta}

\begin{enumerate}[leftmargin=*]
  \item \facil{} Calcula: $a)\; \dfrac{1}{8} + \dfrac{3}{8} \qquad b)\; \dfrac{7}{10} - \dfrac{2}{10} \qquad c)\; \dfrac{4}{9} + \dfrac{5}{9}$

  \item \facil{} Calcula (producto cruzado): $a)\; \dfrac{1}{2} + \dfrac{1}{3} \qquad b)\; \dfrac{2}{5} - \dfrac{1}{4} \qquad c)\; \dfrac{3}{7} + \dfrac{2}{5}$

  \item \facil{} Calcula (mcm): $a)\; \dfrac{1}{6} + \dfrac{3}{4} \qquad b)\; \dfrac{5}{8} - \dfrac{1}{3} \qquad c)\; \dfrac{2}{5} + \dfrac{7}{10}$

  \item \medio{} Calcula: $a)\; 1\dfrac{1}{3} + \dfrac{5}{6} \qquad b)\; 3\dfrac{1}{4} - 1\dfrac{1}{2}$

  \item \medio{} Calcula: $a)\; \dfrac{3}{4} + \dfrac{5}{6} - \dfrac{1}{3} \qquad b)\; \dfrac{7}{12} - \dfrac{1}{4} + \dfrac{2}{3}$

  \item \medio{} Para una receta se necesitan $\dfrac{2}{3}$ de taza de azúcar y $\dfrac{1}{4}$ de taza de harina. \textquestiondown Cuántas tazas de ingredientes se usan en total?

  \item \dificil{} Entre tres amigos se reparte una herencia: el primero recibe $\dfrac{1}{3}$, el segundo $\dfrac{1}{4}$ y el tercero el resto. \textquestiondown Qué fracción recibe el tercero?
\end{enumerate}

\subsection{Multiplicación y división}

\begin{cuadroteoria}
  \textbf{Multiplicación:} se multiplican numeradores entre sí y denominadores entre sí.
  \[
    \frac{a}{b} \cdot \frac{c}{d} = \frac{a \cdot c}{b \cdot d}
  \]

  \textbf{División:} se multiplica la primera fracción por el \textbf{recíproco} de la segunda.
  \[
    \frac{a}{b} \div \frac{c}{d} = \frac{a}{b} \cdot \frac{d}{c} = \frac{a \cdot d}{b \cdot c}
  \]
\end{cuadroteoria}

\begin{cuadroadvertencia}
  \textbf{Cuidado:}
  \begin{itemize}[leftmargin=*]
    \item Para dividir se invierte la \emph{segunda} fracción, no la primera.
    \item Un entero se escribe como fracción con denominador $1$: $3 = \dfrac{3}{1}$.
    \item Antes de multiplicar, \emph{simplifica} factores comunes para trabajar con números más pequeños.
  \end{itemize}
\end{cuadroadvertencia}

\noindent\textbf{Ejemplo resuelto 4:} $\dfrac{3}{8} \cdot \dfrac{4}{9}$. Simplificamos antes de multiplicar: $\dfrac{\cancel{3}}{8} \cdot \dfrac{4}{\cancel{9}} = \dfrac{1}{\cancel{8}^{\,2}} \cdot \dfrac{\cancel{4}^{\,1}}{3} = \dfrac{1}{6}$.

\noindent\textbf{Ejemplo resuelto 5:} $\dfrac{5}{6} \div \dfrac{2}{3} = \dfrac{5}{6} \cdot \dfrac{3}{2} = \dfrac{15}{12} = \dfrac{5}{4} = 1\dfrac{1}{4}$.

\subsection{Ejercicios: Multiplicación y División}

\begin{enumerate}[leftmargin=*]
  \item \facil{} Calcula: $a)\; \dfrac{2}{5} \cdot \dfrac{3}{7} \qquad b)\; \dfrac{1}{4} \cdot \dfrac{8}{9} \qquad c)\; \dfrac{5}{6} \cdot \dfrac{3}{10}$

  \item \facil{} Calcula: $a)\; \dfrac{3}{4} \div \dfrac{2}{5} \qquad b)\; \dfrac{7}{8} \div \dfrac{7}{8} \qquad c)\; \dfrac{9}{10} \div \dfrac{3}{5}$

  \item \facil{} Calcula: $a)\; 5 \cdot \dfrac{2}{3} \qquad b)\; \dfrac{3}{4} \cdot 8 \qquad c)\; \dfrac{1}{2} \div 4$

  \item \medio{} Calcula: $a)\; 2\dfrac{1}{3} \cdot 1\dfrac{1}{2} \qquad b)\; 3\dfrac{3}{4} \div \dfrac{5}{2}$

  \item \medio{} Calcula: $a)\; \dfrac{3}{5} \cdot \dfrac{10}{9} \cdot \dfrac{3}{4} \qquad b)\; \dfrac{2}{3} \div \dfrac{4}{9} \div \dfrac{3}{2}$

  \item \medio{} Un terreno de $\dfrac{3}{4}$ de hectárea se divide en $6$ parcelas iguales. \textquestiondown Qué fracción de hectárea mide cada parcela?

  \item \dificil{} Calcula: $\left(\dfrac{2}{3} \cdot \dfrac{9}{8}\right) \div \left(\dfrac{5}{6} \cdot \dfrac{3}{10}\right)$
\end{enumerate}

\newpage

% ============================================================
\section{Propiedades de la Suma y la Multiplicación en $\Q$}
% ============================================================

\begin{cuadroteoria}
  \textbf{Propiedades de la suma} (con $a, b, c \in \Q$):
  \begin{itemize}[leftmargin=*]
    \item \textbf{Conmutativa:} $a + b = b + a$. \quad Ejemplo: $\frac{1}{3} + \frac{1}{4} = \frac{1}{4} + \frac{1}{3} = \frac{7}{12}$.
    \item \textbf{Asociativa:} $(a + b) + c = a + (b + c)$.
    \item \textbf{Elemento neutro:} $a + 0 = a$. El neutro de la suma es el $0$.
    \item \textbf{Elemento opuesto:} todo $a$ tiene su opuesto $-a$, y $a + (-a) = 0$. \quad Ejemplo: $\frac{3}{5} + \left(-\frac{3}{5}\right) = 0$.
  \end{itemize}
\end{cuadroteoria}

\begin{cuadroteoria}
  \textbf{Propiedades de la multiplicación} (con $a, b, c \in \Q$):
  \begin{itemize}[leftmargin=*]
    \item \textbf{Conmutativa:} $a \cdot b = b \cdot a$.
    \item \textbf{Asociativa:} $(a \cdot b) \cdot c = a \cdot (b \cdot c)$.
    \item \textbf{Elemento neutro:} $a \cdot 1 = a$. El neutro de la multiplicación es el $1$.
    \item \textbf{Elemento inverso:} si $a \neq 0$, su inverso es $\frac{1}{a}$, y $a \cdot \frac{1}{a} = 1$. \quad Ejemplo: $\frac{2}{3} \cdot \frac{3}{2} = 1$.
  \end{itemize}

  \textbf{Propiedad distributiva} (une suma y multiplicación):
  \[
    a \cdot (b + c) = a \cdot b + a \cdot c
  \]
\end{cuadroteoria}

\begin{cuadroadvertencia}
  \textbf{No confundas opuesto con inverso:}
  \begin{itemize}[leftmargin=*]
    \item \textbf{Opuesto} de $\frac{3}{5}$ es $-\frac{3}{5}$ \quad (suma: $\frac{3}{5} + -\frac{3}{5} = 0$).
    \item \textbf{Inverso} de $\frac{3}{5}$ es $\frac{5}{3}$ \quad (multiplica: $\frac{3}{5} \cdot \frac{5}{3} = 1$).
    \item El $0$ \textbf{no tiene inverso} (no se puede dividir entre cero).
  \end{itemize}
\end{cuadroadvertencia}

\noindent\textbf{Ejemplo resuelto 1:} Identifica la propiedad en $\dfrac{3}{4} \cdot \left(\dfrac{1}{2} + \dfrac{1}{3}\right) = \dfrac{3}{4} \cdot \dfrac{1}{2} + \dfrac{3}{4} \cdot \dfrac{1}{3}$.

Es la \textbf{propiedad distributiva}. Comprobamos: el lado derecho es $\dfrac{3}{8} + \dfrac{1}{4} = \dfrac{3}{8} + \dfrac{2}{8} = \dfrac{5}{8}$, y el izquierdo es $\dfrac{3}{4} \cdot \dfrac{5}{6} = \dfrac{15}{24} = \dfrac{5}{8}$. \checkmark

\subsection{Ejercicios: Propiedades}

\begin{enumerate}[leftmargin=*]
  \item \facil{} Nombra la propiedad aplicada:
        \[ a)\; \frac{2}{5} + \frac{3}{7} = \frac{3}{7} + \frac{2}{5} \quad b)\; \frac{3}{4} \cdot 1 = \frac{3}{4} \quad c)\; \frac{1}{2} + 0 = \frac{1}{2} \]

  \item \facil{} Escribe el \textbf{opuesto} y el \textbf{inverso} de: $a)\; \dfrac{3}{5} \qquad b)\; -\dfrac{2}{7} \qquad c)\; 4$

  \item \facil{} Nombra la propiedad: $\left(\dfrac{1}{2} + \dfrac{1}{3}\right) + \dfrac{1}{6} = \dfrac{1}{2} + \left(\dfrac{1}{3} + \dfrac{1}{6}\right)$.

  \item \medio{} Aplica la propiedad distributiva y calcula: $\dfrac{2}{3} \cdot \left(\dfrac{1}{4} + \dfrac{1}{2}\right)$.

  \item \medio{} Comprueba la propiedad asociativa de la multiplicación con $\dfrac{1}{2}$, $\dfrac{2}{3}$ y $\dfrac{3}{4}$.

  \item \medio{} \textquestiondown Cuál es el inverso de $\dfrac{5}{4}$? \textquestiondown Y su opuesto? \textquestiondown Qué obtienes al sumarlos?

  \item \dificil{} Usa la propiedad distributiva para calcular de forma rápida: $\left(\dfrac{5}{7} + \dfrac{2}{7}\right) \cdot \dfrac{3}{4}$.
\end{enumerate}

% ============================================================
\section{Operaciones Combinadas}
% ============================================================

\subsection{Jerarquía de las operaciones}

En una expresión con varias operaciones \textbf{no} se resuelve de izquierda a derecha: existe una \textbf{jerarquía} que se debe respetar.

\begin{cuadroteoria}
  \textbf{Orden de resolución:}
  \begin{enumerate}[leftmargin=*]
    \item \textbf{Signos de agrupación}, de adentro hacia afuera: paréntesis $(\ )$, corchetes $[\ ]$, llaves $\{\ \}$.
    \item \textbf{Multiplicaciones y divisiones}, de izquierda a derecha.
    \item \textbf{Sumas y restas}, de izquierda a derecha.
    \item \textbf{Simplifica} el resultado final.
  \end{enumerate}
\end{cuadroteoria}

\noindent\textbf{Ejemplo resuelto 1 (sin agrupación):} $\dfrac{1}{2} + \dfrac{2}{3} \cdot \dfrac{3}{4}$.

Primero la multiplicación: $\dfrac{2}{3} \cdot \dfrac{3}{4} = \dfrac{6}{12} = \dfrac{1}{2}$. Luego la suma: $\dfrac{1}{2} + \dfrac{1}{2} = 1$.

\noindent\textbf{Ejemplo resuelto 2 (con agrupación):} $\left[\left(\dfrac{1}{3} + \dfrac{1}{6}\right) \cdot 2 - \dfrac{1}{4}\right] \div \dfrac{5}{8}$.

Paréntesis: $\dfrac{1}{3} + \dfrac{1}{6} = \dfrac{3}{6} = \dfrac{1}{2}$. Multiplicación: $\dfrac{1}{2} \cdot 2 = 1$. Resta: $1 - \dfrac{1}{4} = \dfrac{3}{4}$. División final: $\dfrac{3}{4} \div \dfrac{5}{8} = \dfrac{3}{4} \cdot \dfrac{8}{5} = \dfrac{6}{5}$.

\begin{cuadroadvertencia}
  \textbf{\textquestiondown Por qué importa el orden?} $\dfrac{1}{2} + \dfrac{1}{3} \cdot \dfrac{3}{4} = 1$, pero $\left(\dfrac{1}{2} + \dfrac{1}{3}\right) \cdot \dfrac{3}{4} = \dfrac{5}{6} \cdot \dfrac{3}{4} = \dfrac{5}{8}$. Un paréntesis cambia todo el resultado.
\end{cuadroadvertencia}

\subsection{Ejercicios: Operaciones Combinadas}

\begin{enumerate}[leftmargin=*]
  \item \facil{} Calcula: $a)\; \dfrac{1}{2} \cdot \dfrac{4}{5} + \dfrac{1}{5} \qquad b)\; \dfrac{3}{4} - \dfrac{1}{2} \cdot \dfrac{1}{2}$

  \item \facil{} Calcula: $a)\; \dfrac{5}{6} \div \dfrac{1}{3} - \dfrac{1}{2} \qquad b)\; \dfrac{2}{3} + \dfrac{3}{4} \div \dfrac{9}{8}$

  \item \medio{} Calcula: $a)\; \left(\dfrac{1}{4} + \dfrac{3}{4}\right) \cdot \dfrac{2}{5} \qquad b)\; \left(\dfrac{2}{3} - \dfrac{1}{6}\right) \div \dfrac{1}{2}$

  \item \medio{} Calcula: $a)\; \left(2\dfrac{1}{2} - \dfrac{3}{4}\right) \cdot \dfrac{4}{7} \qquad b)\; \dfrac{3}{5} \cdot \left(\dfrac{5}{6} + \dfrac{1}{3}\right)$

  \item \medio{} Calcula: $\dfrac{\frac{2}{3} + \frac{1}{6}}{\frac{5}{6} - \frac{1}{2}}$

  \item \medio{} Calcula: $\left[\left(\dfrac{1}{2} + \dfrac{1}{4}\right) \cdot \dfrac{4}{3} + \dfrac{1}{6}\right] \div \dfrac{5}{12}$

  \item \dificil{} Calcula: $\dfrac{\frac{3}{4} \cdot \frac{8}{9} - \frac{1}{3}}{\frac{5}{6} \div \frac{5}{12} + \frac{1}{2}}$

  \item \dificil{} Plantea y resuelve: tengo $\dfrac{3}{4}$ de litro de jugo. Bebo $\dfrac{1}{3}$ del total y luego regalo la mitad de lo que queda. \textquestiondown Cuánto jugo me queda?
\end{enumerate}

\newpage

% ============================================================
\section{Potenciación en $\Q$}
% ============================================================

\subsection{Definición y propiedades}

\begin{cuadroteoria}
  \textbf{Potencia de una fracción:} se elevan numerador y denominador al exponente.
  \[
    \left(\frac{a}{b}\right)^n = \frac{a^n}{b^n} \qquad (b \neq 0)
  \]

  \textbf{Signo del resultado:}
  \begin{itemize}[leftmargin=*]
    \item Exponente \textbf{par}: el resultado es \textbf{positivo}. \quad $\left(-\dfrac{2}{3}\right)^2 = \dfrac{4}{9}$.
    \item Exponente \textbf{impar}: se conserva el signo de la base. \quad $\left(-\dfrac{2}{3}\right)^3 = -\dfrac{8}{27}$.
  \end{itemize}
\end{cuadroteoria}

\begin{cuadroteoria}
  \textbf{Propiedades de las potencias} (con $a, b \neq 0$):
  \begin{center}
    \renewcommand{\arraystretch}{1.4}
    \begin{tabular}{ll}
      \toprule
      \textbf{Propiedad}       & \textbf{Fórmula}                                 \\
      \midrule
      Producto de potencias    & $a^m \cdot a^n = a^{m+n}$                        \\
      Cociente de potencias    & $\dfrac{a^m}{a^n} = a^{m-n}$                     \\
      Potencia de una potencia & $(a^m)^n = a^{m \cdot n}$                        \\
      Potencia de un producto  & $(a \cdot b)^n = a^n \cdot b^n$                  \\
      Potencia de un cociente  & $\left(\dfrac{a}{b}\right)^n = \dfrac{a^n}{b^n}$ \\
      Exponente cero           & $a^0 = 1$                                        \\
      Exponente negativo       & $a^{-n} = \dfrac{1}{a^n}$ \quad y \quad $\left(\dfrac{a}{b}\right)^{-n} = \left(\dfrac{b}{a}\right)^{n}$ \\
      \bottomrule
    \end{tabular}
  \end{center}
\end{cuadroteoria}

\begin{cuadroadvertencia}
  \textbf{Errores clásicos:}
  \begin{itemize}[leftmargin=*]
    \item Un exponente negativo \textbf{invierte} la fracción: $\left(\dfrac{2}{5}\right)^{-2} = \left(\dfrac{5}{2}\right)^2 = \dfrac{25}{4}$.
    \item En el cociente los exponentes se \textbf{restan}: $\dfrac{a^m}{a^n} = a^{m-n}$ (no se multiplican).
    \item La potencia afecta a \emph{toda} la base: $\left(-\dfrac{2}{3}\right)^2 = \dfrac{4}{9}$, pero $-\dfrac{2^2}{3} = -\dfrac{4}{3}$.
  \end{itemize}
\end{cuadroadvertencia}

\noindent\textbf{Ejemplo resuelto 1:} $\left(-\dfrac{3}{2}\right)^{-3} = \left(-\dfrac{2}{3}\right)^3 = \dfrac{(-2)^3}{3^3} = -\dfrac{8}{27}$.

\noindent\textbf{Ejemplo resuelto 2 (simplificación tipo examen):} Simplifica $\left(\dfrac{81a^3b^2}{27a^5b^6}\right)^{-2}$.

\begin{cuadrosolucion}
  \textbf{Paso 1:} simplificamos la fracción interior. $\dfrac{81}{27} = 3$ y aplicamos cociente de potencias a cada variable:
  \[
    \frac{81a^3b^2}{27a^5b^6} = 3\,a^{3-5}b^{2-6} = 3a^{-2}b^{-4}
  \]
  \textbf{Paso 2:} aplicamos el exponente $-2$ (potencia de un producto y de una potencia):
  \[
    \left(3a^{-2}b^{-4}\right)^{-2} = 3^{-2}\,a^{(-2)(-2)}\,b^{(-4)(-2)} = 3^{-2}a^{4}b^{8}
  \]
  \textbf{Paso 3:} escribimos sin exponentes negativos, $3^{-2} = \dfrac{1}{9}$:
  \[
    \boxed{\ \left(\frac{81a^3b^2}{27a^5b^6}\right)^{-2} = \frac{a^{4}b^{8}}{9}\ }
  \]
\end{cuadrosolucion}

\subsection{Ejercicios: Potenciación en $\Q$}

\begin{enumerate}[leftmargin=*]
  \item \facil{} Calcula: $a)\; \left(\dfrac{3}{4}\right)^2 \qquad b)\; \left(\dfrac{2}{5}\right)^3 \qquad c)\; \left(-\dfrac{1}{3}\right)^2$

  \item \facil{} Calcula: $a)\; \left(-\dfrac{2}{3}\right)^3 \qquad b)\; \left(\dfrac{4}{5}\right)^0 \qquad c)\; \left(-\dfrac{1}{2}\right)^5$

  \item \facil{} Calcula: $a)\; \left(\dfrac{3}{5}\right)^{-1} \qquad b)\; \left(\dfrac{2}{3}\right)^{-2} \qquad c)\; \left(-\dfrac{5}{4}\right)^{-2}$

  \item \medio{} Aplica propiedades: $a)\; \left(\dfrac{2}{7}\right)^3 \cdot \left(\dfrac{2}{7}\right)^4 \qquad b)\; \dfrac{\left(\frac{3}{5}\right)^7}{\left(\frac{3}{5}\right)^4}$

  \item \medio{} Simplifica: $a)\; \left(\dfrac{4x^3}{y^2}\right)^2 \qquad b)\; \left(\dfrac{2a^2b}{3c}\right)^3$

  \item \medio{} Simplifica: $\left(\dfrac{16x^4y^3}{8x^6y^7}\right)^{-1}$

  \item \medio{} Simplifica: $\left(\dfrac{27a^5b^2}{9a^2b^6}\right)^{-2}$

  \item \dificil{} Simplifica: $\left(\dfrac{2x^2y^3}{z^2}\right)^{-3} \cdot \left(\dfrac{x^4z^6}{8y^9}\right)$

  \item \dificil{} El volumen de un cubo es $\left(\dfrac{2}{3}\right)^3\ \text{m}^3$. \textquestiondown Cuál es la medida de su arista?
\end{enumerate}

\newpage

% ============================================================
\section{Ecuaciones Lineales en $\Q$}
% ============================================================

\subsection{\textquestiondown Cómo se resuelven?}

Una \textbf{ecuación lineal} en $\Q$ es una igualdad donde la incógnita aparece a la primera potencia y puede haber fracciones. Conviene \textbf{eliminar los denominadores} multiplicando \emph{toda} la ecuación por el mcm.

\begin{cuadropasos}
  \textbf{Método de resolución:}
  \begin{enumerate}[leftmargin=*]
    \item \textbf{Eliminar denominadores:} multiplica \emph{toda} la ecuación por el mcm de los denominadores.
    \item \textbf{Distribuir:} aplica la propiedad distributiva en los paréntesis.
    \item \textbf{Reducir:} agrupa los términos con $x$ en un lado y los números en el otro.
    \item \textbf{Despejar:} divide por el coeficiente de $x$.
    \item \textbf{Comprobar:} sustituye la solución en la ecuación original.
  \end{enumerate}
\end{cuadropasos}

\noindent\textbf{Ejemplo resuelto 1:} Resuelve $\dfrac{x}{3} + \dfrac{x}{2} = 5$.

mcm$(3,2) = 6$. Multiplicamos toda la ecuación por $6$: $2x + 3x = 30 \;\Rightarrow\; 5x = 30 \;\Rightarrow\; x = 6$. Comprobación: $\dfrac{6}{3} + \dfrac{6}{2} = 2 + 3 = 5$. \checkmark

\noindent\textbf{Ejemplo resuelto 2:} Resuelve $\dfrac{2x - 1}{3} = \dfrac{x + 2}{2}$.

mcm$(3,2) = 6$: $2(2x - 1) = 3(x + 2) \;\Rightarrow\; 4x - 2 = 3x + 6 \;\Rightarrow\; x = 8$. \quad Comprobación: $\dfrac{15}{3} = 5$ y $\dfrac{10}{2} = 5$. \checkmark

\begin{cuadroadvertencia}
  \textbf{Error clásico:} multiplicar solo \emph{algunos} términos por el mcm. El mcm debe multiplicar \textbf{toda} la ecuación, incluidos los términos enteros. Y al distribuir un signo menos: $-\,(x - 2) = -x + 2$.
\end{cuadroadvertencia}

\subsection{Ejercicios: Ecuaciones}

\begin{enumerate}[leftmargin=*]
  \item \facil{} Resuelve: $\dfrac{x}{5} = 3$

  \item \facil{} Resuelve: $\dfrac{x}{4} + 2 = 5$

  \item \facil{} Resuelve: $\dfrac{x}{2} + \dfrac{x}{3} = 10$

  \item \facil{} Resuelve: $\dfrac{2x}{5} = \dfrac{6}{5}$

  \item \medio{} Resuelve: $\dfrac{x + 1}{3} = \dfrac{x - 2}{4}$

  \item \medio{} Resuelve: $\dfrac{3x}{4} - \dfrac{x}{8} = 5$

  \item \medio{} Resuelve: $\dfrac{2x - 3}{5} + \dfrac{x}{2} = 3$

  \item \medio{} Resuelve: $\dfrac{x}{3} + \dfrac{x + 1}{2} = \dfrac{7}{6}$

  \item \medio{} Resuelve: $\dfrac{2x + 1}{3} - \dfrac{x - 2}{4} = \dfrac{5x}{6}$

  \item \dificil{} Resuelve: $\dfrac{x - 1}{2} - \dfrac{x - 2}{3} = \dfrac{x - 3}{4}$

  \item \dificil{} La cuarta parte de un número aumentada en $\dfrac{1}{2}$ es igual a la sexta parte del número más $2$. \textquestiondown Cuál es el número?

  \item \dificil{} En una clase, $\dfrac{2}{5}$ de los estudiantes practican fútbol, $\dfrac{1}{3}$ practican básquetbol y los $8$ restantes practican otro deporte. \textquestiondown Cuántos estudiantes tiene la clase?
\end{enumerate}

\newpage

% ============================================================
\ifmwclaves
  \section{Clave de Respuestas}
  % ============================================================

  \subsection*{Fracciones y Conversiones}
  \begin{enumerate}[leftmargin=*, label=\textbf{\arabic*.}]
    \item $a)$ propia; \quad $b)$ impropia; \quad $c)$ impropia ($=1$); \quad $d)$ propia.
    \item $a)\; \dfrac{11}{4}$; \quad $b)\; \dfrac{31}{6}$; \quad $c)\; \dfrac{26}{3}$.
    \item $a)\; 3\dfrac{2}{5}$; \quad $b)\; 3\dfrac{1}{7}$; \quad $c)\; 7\dfrac{3}{4}$.
    \item $a)\; \dfrac{2}{3}$; \quad $b)\; \dfrac{2}{3}$; \quad $c)\; \dfrac{2}{3}$.
    \item $a)\; \dfrac{6}{15}$; \quad $b)\; \dfrac{28}{32}$; \quad $c)\; \dfrac{20}{24}$.
    \item $\dfrac{7}{12} < \dfrac{3}{4} < \dfrac{5}{6}$.
    \item Quedan $\dfrac{5}{8}$ (propia).
    \item $2 \cdot \dfrac{3}{4} = \dfrac{3}{2} = 1\dfrac{1}{2}$ km.
    \item $5\dfrac{3}{8} = \dfrac{43}{8}$, y $\dfrac{43}{8} = 5\dfrac{3}{8}$. \checkmark
  \end{enumerate}

  \subsection*{Suma y Resta}
  \begin{enumerate}[leftmargin=*, label=\textbf{\arabic*.}]
    \item $a)\; \dfrac{1}{2}$; \quad $b)\; \dfrac{1}{2}$; \quad $c)\; 1$.
    \item $a)\; \dfrac{5}{6}$; \quad $b)\; \dfrac{3}{20}$; \quad $c)\; \dfrac{29}{35}$.
    \item $a)\; \dfrac{11}{12}$; \quad $b)\; \dfrac{7}{24}$; \quad $c)\; \dfrac{11}{10}$.
    \item $a)\; 1\dfrac{1}{3} + \dfrac{5}{6} = \dfrac{4}{3} + \dfrac{5}{6} = \dfrac{13}{6} = 2\dfrac{1}{6}$; \quad $b)\; \dfrac{13}{4} - \dfrac{3}{2} = \dfrac{13}{4} - \dfrac{6}{4} = \dfrac{7}{4} = 1\dfrac{3}{4}$.
    \item $a)\; \dfrac{9 + 10 - 4}{12} = \dfrac{15}{12} = \dfrac{5}{4}$; \quad $b)\; \dfrac{7 - 3 + 8}{12} = 1$.
    \item $\dfrac{2}{3} + \dfrac{1}{4} = \dfrac{11}{12}$ de taza.
    \item $1 - \dfrac{1}{3} - \dfrac{1}{4} = 1 - \dfrac{7}{12} = \dfrac{5}{12}$.
  \end{enumerate}

  \subsection*{Multiplicación y División}
  \begin{enumerate}[leftmargin=*, label=\textbf{\arabic*.}]
    \item $a)\; \dfrac{6}{35}$; \quad $b)\; \dfrac{2}{9}$; \quad $c)\; \dfrac{1}{4}$.
    \item $a)\; \dfrac{15}{8}$; \quad $b)\; 1$; \quad $c)\; \dfrac{3}{2}$.
    \item $a)\; \dfrac{10}{3}$; \quad $b)\; 6$; \quad $c)\; \dfrac{1}{8}$.
    \item $a)\; \dfrac{7}{3} \cdot \dfrac{3}{2} = \dfrac{7}{2} = 3\dfrac{1}{2}$; \quad $b)\; \dfrac{15}{4} \div \dfrac{5}{2} = \dfrac{15}{4} \cdot \dfrac{2}{5} = \dfrac{3}{2} = 1\dfrac{1}{2}$.
    \item $a)\; \dfrac{3 \cdot 10 \cdot 3}{5 \cdot 9 \cdot 4} = \dfrac{90}{180} = \dfrac{1}{2}$; \quad $b)\; \dfrac{2}{3} \cdot \dfrac{9}{4} \cdot \dfrac{2}{3} = \dfrac{36}{36} = 1$.
    \item $\dfrac{3}{4} \div 6 = \dfrac{3}{24} = \dfrac{1}{8}$ de hectárea.
    \item $\left(\dfrac{2}{3} \cdot \dfrac{9}{8}\right) = \dfrac{3}{4}$ y $\left(\dfrac{5}{6} \cdot \dfrac{3}{10}\right) = \dfrac{1}{4}$; entonces $\dfrac{3}{4} \div \dfrac{1}{4} = 3$.
  \end{enumerate}

  \subsection*{Propiedades}
  \begin{enumerate}[leftmargin=*, label=\textbf{\arabic*.}]
    \item $a)$ conmutativa de la suma; \quad $b)$ elemento neutro de la multiplicación; \quad $c)$ elemento neutro de la suma.
    \item $a)$ opuesto $-\dfrac{3}{5}$, inverso $\dfrac{5}{3}$; \quad $b)$ opuesto $\dfrac{2}{7}$, inverso $-\dfrac{7}{2}$; \quad $c)$ opuesto $-4$, inverso $\dfrac{1}{4}$.
    \item Propiedad asociativa de la suma.
    \item $\dfrac{2}{3} \cdot \dfrac{1}{4} + \dfrac{2}{3} \cdot \dfrac{1}{2} = \dfrac{1}{6} + \dfrac{1}{3} = \dfrac{1}{2}$.
    \item $\left(\dfrac{1}{2} \cdot \dfrac{2}{3}\right) \cdot \dfrac{3}{4} = \dfrac{1}{3} \cdot \dfrac{3}{4} = \dfrac{1}{4}$ y $\dfrac{1}{2} \cdot \left(\dfrac{2}{3} \cdot \dfrac{3}{4}\right) = \dfrac{1}{2} \cdot \dfrac{1}{2} = \dfrac{1}{4}$. \checkmark (se cumple)
    \item Inverso $\dfrac{4}{5}$, opuesto $-\dfrac{5}{4}$; su suma es $-\dfrac{9}{20}$.
    \item $\left(\dfrac{5}{7} + \dfrac{2}{7}\right) \cdot \dfrac{3}{4} = 1 \cdot \dfrac{3}{4} = \dfrac{3}{4}$.
  \end{enumerate}

  \subsection*{Operaciones Combinadas}
  \begin{enumerate}[leftmargin=*, label=\textbf{\arabic*.}]
    \item $a)\; \dfrac{2}{5} + \dfrac{1}{5} = \dfrac{3}{5}$; \quad $b)\; \dfrac{3}{4} - \dfrac{1}{4} = \dfrac{1}{2}$.
    \item $a)\; \dfrac{5}{2} - \dfrac{1}{2} = 2$; \quad $b)\; \dfrac{2}{3} + \dfrac{2}{3} = \dfrac{4}{3}$.
    \item $a)\; 1 \cdot \dfrac{2}{5} = \dfrac{2}{5}$; \quad $b)\; \dfrac{1}{2} \div \dfrac{1}{2} = 1$.
    \item $a)\; \dfrac{7}{4} \cdot \dfrac{4}{7} = 1$; \quad $b)\; \dfrac{3}{5} \cdot \dfrac{7}{6} = \dfrac{7}{10}$.
    \item $\dfrac{\frac{5}{6}}{\frac{1}{3}} = \dfrac{5}{6} \cdot 3 = \dfrac{5}{2}$.
    \item $\left[\dfrac{3}{4} \cdot \dfrac{4}{3} + \dfrac{1}{6}\right] \div \dfrac{5}{12} = \left[1 + \dfrac{1}{6}\right] \div \dfrac{5}{12} = \dfrac{7}{6} \cdot \dfrac{12}{5} = \dfrac{14}{5}$.
    \item $\dfrac{\frac{2}{3} - \frac{1}{3}}{\frac{2}{1} + \frac{1}{2}} = \dfrac{\frac{1}{3}}{\frac{5}{2}} = \dfrac{1}{3} \cdot \dfrac{2}{5} = \dfrac{2}{15}$.
    \item Quedan: $\dfrac{3}{4} - \dfrac{1}{4} = \dfrac{1}{2}$ L bebo; queda $\dfrac{1}{2}$; regalo la mitad, queda $\dfrac{1}{4}$ L.
  \end{enumerate}

  \subsection*{Potenciación en $\Q$}
  \begin{enumerate}[leftmargin=*, label=\textbf{\arabic*.}]
    \item $a)\; \dfrac{9}{16}$; \quad $b)\; \dfrac{8}{125}$; \quad $c)\; \dfrac{1}{9}$.
    \item $a)\; -\dfrac{8}{27}$; \quad $b)\; 1$; \quad $c)\; -\dfrac{1}{32}$.
    \item $a)\; \dfrac{5}{3}$; \quad $b)\; \dfrac{9}{4}$; \quad $c)\; \dfrac{16}{25}$.
    \item $a)\; \left(\dfrac{2}{7}\right)^7 = \dfrac{128}{823\,543}$; \quad $b)\; \left(\dfrac{3}{5}\right)^3 = \dfrac{27}{125}$.
    \item $a)\; \dfrac{16x^6}{y^4}$; \quad $b)\; \dfrac{8a^6b^3}{27c^3}$.
    \item $\left(\dfrac{16x^4y^3}{8x^6y^7}\right)^{-1} = \left(2x^{-2}y^{-4}\right)^{-1} = 2^{-1}x^{2}y^{4} = \dfrac{x^2y^4}{2}$.
    \item $\left(\dfrac{27a^5b^2}{9a^2b^6}\right)^{-2} = \left(3a^3b^{-4}\right)^{-2} = 3^{-2}a^{-6}b^{8} = \dfrac{b^8}{9a^6}$.
    \item $\left(\dfrac{2x^2y^3}{z^2}\right)^{-3} = \dfrac{z^6}{8x^6y^9}$; multiplicando por $\dfrac{x^4z^6}{8y^9}$ se obtiene $\dfrac{x^4z^{12}}{64x^6y^{18}} = \dfrac{z^{12}}{64x^2y^{18}}$.
    \item La arista mide $\dfrac{2}{3}$ m.
  \end{enumerate}

  \subsection*{Ecuaciones Lineales}
  \begin{enumerate}[leftmargin=*, label=\textbf{\arabic*.}]
    \item $x = 15$.
    \item $\dfrac{x}{4} = 3 \Rightarrow x = 12$.
    \item $\dfrac{5x}{6} = 10 \Rightarrow x = 12$.
    \item $2x = 6 \Rightarrow x = 3$.
    \item $4(x + 1) = 3(x - 2) \Rightarrow 4x + 4 = 3x - 6 \Rightarrow x = -10$.
    \item $6x - x = 40 \Rightarrow 5x = 40 \Rightarrow x = 8$.
    \item $2(2x - 3) + 5x = 30 \Rightarrow 9x - 6 = 30 \Rightarrow x = 4$.
    \item $2x + 3(x + 1) = 7 \Rightarrow 5x + 3 = 7 \Rightarrow x = \dfrac{4}{5}$.
    \item $4(2x + 1) - 3(x - 2) = 10x \Rightarrow 5x + 10 = 10x \Rightarrow x = 2$.
    \item $6(x - 1) - 4(x - 2) = 3(x - 3) \Rightarrow 2x + 2 = 3x - 9 \Rightarrow x = 11$.
    \item $\dfrac{x}{4} + \dfrac{1}{2} = \dfrac{x}{6} + 2 \Rightarrow 3x + 6 = 2x + 24 \Rightarrow x = 18$.
    \item $x - \dfrac{2x}{5} - \dfrac{x}{3} = 8 \Rightarrow \dfrac{4x}{15} = 8 \Rightarrow x = 30$ estudiantes.
  \end{enumerate}

  \vspace{10pt}
  \begin{cuadrosolucion}
    \textbf{\textquestiondown Cómo te fue?} Si resolviste hasta el final, ya dominas todo el Tema 1 de Números Racionales. Comparar tus resultados con esta clave es parte del aprendizaje: cada error es una pista para mejorar. ¡Nos vemos en el Tema 2!
  \end{cuadrosolucion}

\fi
\end{document}
